\documentclass[11pt,a4paper]{article}
\usepackage[T1]{fontenc}
\usepackage[utf8]{inputenc}
\usepackage{lmodern}
\usepackage{microtype}
\usepackage[a4paper,margin=2.25cm]{geometry}
\usepackage{longtable,booktabs,array,calc}
\usepackage{fancyvrb}
\usepackage{xcolor}
\usepackage{hyperref}
\usepackage{xurl}
\usepackage{fancyhdr}
\usepackage{enumitem}
\usepackage{parskip}
\usepackage{titlesec}
\definecolor{ddtadark}{RGB}{30,38,48}
\definecolor{ddtablue}{RGB}{38,82,126}
\hypersetup{colorlinks=true,linkcolor=ddtablue,urlcolor=ddtablue,citecolor=ddtablue}
\pagestyle{fancy}
\fancyhf{}
\lhead{Documentation-Driven Threat Analysis}
\rhead{BA3-T3 pressure test}
\cfoot{\thepage}
\setlength{\headheight}{14pt}
\setlist[itemize]{leftmargin=1.5em,itemsep=0.2em,topsep=0.25em}
\setlist[enumerate]{leftmargin=1.7em,itemsep=0.2em,topsep=0.25em}
\titleformat{\section}{\Large\bfseries\color{ddtadark}}{}{0em}{}
\titleformat{\subsection}{\large\bfseries\color{ddtadark}}{}{0em}{}
\titleformat{\subsubsection}{\normalsize\bfseries\color{ddtadark}}{}{0em}{}
\providecommand{\tightlist}{\setlength{\itemsep}{0pt}\setlength{\parskip}{0pt}}
\setlength{\emergencystretch}{4em}
\hbadness=10000
\hfuzz=20pt
\sloppy
\begin{document}
\section{DDTA BA3-T3 derivation-rule reproducibility and change-impact
lineage pressure test -
R1}\label{ddta-ba3-t3-derivation-rule-reproducibility-and-change-impact-lineage-pressure-test---r1}

\textbf{Status:} COMPLETED / PROVISIONAL PASS WITH DERIVATION-IMPACT
REFINEMENT

\textbf{Repository baseline tested:}
\texttt{5fc0b92809ece193deaba4206488d78f981f7855}

\textbf{Scope:} only BA3-T3. BA4, formal threat-method overlays,
AnalysisRecord, Common Finding and implementation design remain out of
scope.

\subsection{1. Test objective}\label{test-objective}

BA3-T3 tests three falsifiable propositions left open by T2:

\begin{enumerate}
\def\labelenumi{\arabic{enumi}.}
\tightlist
\item
  \texttt{derivationBasis\ +\ derivationRuleRef} is enough for
  reproducible review only if the basis and rule semantics are
  sufficiently explicit;
\item
  direct source provenance plus BA structure is enough for localized
  change impact only if all meaning-relevant governed context is
  represented;
\item
  the analysis-to-correction feedback loop can remain outside BA source
  authority while still being traceable back into the next governed
  baseline.
\end{enumerate}

The test uses concrete mutation/provider cases rather than constructing
a generic dependency engine.

\subsection{2. Starting lower bound}\label{starting-lower-bound}

Before T3, BA3 had:

\begin{verbatim}
BAOriginAttachment
|- targetElement
|- governedBaselineKey
|- originState
|- sourceLink
|- derivationBasis
`- derivationRuleRef

BABaselineReview
|- targetElement
|- evaluatedBaselineKey
|- reviewState
`- freshness

BACrossBaselineContinuity
|- priorElement
|- targetBaselineKey
|- disposition
|- successorElement
`- continuityBasis
\end{verbatim}

T2 also identified a staleness candidate when known effective governed
context changes, but the exact representation of that context remained
open.

\subsection{3. Falsification attempt A - plain derivation-basis
list}\label{falsification-attempt-a---plain-derivation-basis-list}

\subsubsection{Hypothesis A0}\label{hypothesis-a0}

\begin{quote}
A derived BA element can remain reproducible with an unordered/untyped
list of basis references plus an opaque rule reference.
\end{quote}

\subsubsection{Pressure case}\label{pressure-case}

The order/payment corpus distinguishes provider-specific raw status from
governed project status. A normalization derivation needs at least two
conceptually different inputs:

\begin{verbatim}
provider vocabulary/status evidence
project-governed mapping/contract
\end{verbatim}

If both are stored merely as:

\begin{verbatim}
derivationBasis = [source-A, source-B]
\end{verbatim}

a reviewer cannot determine which source is the mapping authority, which
one provides the raw vocabulary, or whether a tool silently embedded the
mapping.

\subsubsection{Result}\label{result}

\begin{verbatim}
A0 REJECTED
\end{verbatim}

The basis must be role-bound under the referenced rule:

\begin{verbatim}
inputRoleKey -> basisRef
\end{verbatim}

The role key belongs to the rule contract, not BA2 proposition
participation.

\subsection{4. Falsification attempt B - opaque/unversioned derivation
rule}\label{falsification-attempt-b---opaqueunversioned-derivation-rule}

\subsubsection{Hypothesis B0}\label{hypothesis-b0}

\begin{quote}
\texttt{derivationRuleRef\ =\ provider-normalization} is sufficient even
if the rule definition can change in place.
\end{quote}

\subsubsection{Counterexample}\label{counterexample}

Suppose revision R1 maps provider outcomes using one governed status
interpretation and revision R2 changes applicability/output semantics.
If historical derived elements reference only the mutable name,
replaying B0 later can silently use R2 and produce a different BA
result.

\subsubsection{Result}\label{result-1}

\begin{verbatim}
B0 REJECTED
\end{verbatim}

A derivation reference must resolve to an immutable rule revision. The
registry entry must expose input roles, applicability, output semantic
contract and a method-neutral normative definition/rationale.

An implementation can store these as one document or several fields; the
semantic responsibility is what matters.

\subsection{5. Reproducibility
threshold}\label{reproducibility-threshold}

The test deliberately rejects machine-execution as the threshold.

Required:

\begin{verbatim}
same baseline
+ same role-bound basis
+ same immutable rule revision
+ same accepted semantic registries
    -> materially equivalent semantic result
       or an explicit non-applicability/diagnostic outcome
\end{verbatim}

Not required:

\begin{verbatim}
same UUID
same JSON byte order
same graph traversal order
same software implementation
same LLM/tool version
\end{verbatim}

If hidden analyst/tool judgment is required to choose among materially
different outputs, the derivation is not reproducible enough for silent
\texttt{ACCEPTED} status.

\subsection{6. Provider-state normalization
replay}\label{provider-state-normalization-replay}

\subsubsection{Governed payment case}\label{governed-payment-case}

The corpus governs \texttt{PaymentAuthorizationResult} as:

\begin{verbatim}
authorized
 declined
 indeterminate
\end{verbatim}

while provider-specific raw states remain behind the adapter/contract
boundary.

A valid derived normalization needs a governed source that authorizes
the concrete mapping. The method-neutral rule may say:

\begin{quote}
apply the governed provider-to-project status mapping and materialize
the corresponding controlled project semantic result.
\end{quote}

It may not say, without governed evidence:

\begin{verbatim}
provider APPROVED means authorized
\end{verbatim}

\subsubsection{Cross-domain control}\label{cross-domain-control}

The same pattern appears in:

\begin{itemize}
\tightlist
\item
  external WMS reservation results;
\item
  carrier acceptance/status normalization;
\item
  any provider-specific raw result mapped behind a governed project
  contract.
\end{itemize}

\subsubsection{Result}\label{result-2}

\begin{verbatim}
PASS
\end{verbatim}

One rule-registry contract works across the distinct domains without
importing payment/WMS/carrier vocabulary into Base Analysis core
semantics.

\subsection{7. Falsification attempt C - direct provenance is enough for
M4}\label{falsification-attempt-c---direct-provenance-is-enough-for-m4}

\subsubsection{Hypothesis C0}\label{hypothesis-c0}

\begin{quote}
A grounded proposition needs only its direct source links; sibling
governed context can be ignored until the direct FR changes.
\end{quote}

\subsubsection{B0}\label{b0}

The camera corpus has remote recognition and \texttt{FR-3.4} delivers
\texttt{RecognitionCapture} to \texttt{RecognitionProcessor}.

A grounded BA proposition is:

\begin{verbatim}
P0 = transfer(CameraSubsystem,
              RecognitionProcessor,
              RecognitionCapture)
\end{verbatim}

\subsubsection{B1 mutation}\label{b1-mutation}

A sibling representation Decision changes the boundary representation:

\begin{verbatim}
RecognitionCapture -> OpaqueEvidenceReference
\end{verbatim}

while, as an adversarial intermediate state, the old FR text has not yet
been corrected.

With only \texttt{sourceLink(FR-3.4)}, P0 receives no change trigger and
can be falsely treated as current.

\subsubsection{Result}\label{result-3}

\begin{verbatim}
C0 REJECTED
\end{verbatim}

The proposition requires an explicit \texttt{revalidationContext} link
to the representation commitment when that commitment is known to
determine validity/applicability.

\subsection{8. M4 authority test - stale/diagnose, do not
invent}\label{m4-authority-test---stalediagnose-do-not-invent}

Once the sibling representation Decision changes, T3 tests two possible
reactions.

\subsubsection{Alternative 1 - silent
rewrite}\label{alternative-1---silent-rewrite}

\begin{verbatim}
P0 capture-transfer
    -> automatically rewrite to
P1 opaque-reference-transfer
\end{verbatim}

without governed FR correction.

\begin{verbatim}
REJECTED
\end{verbatim}

This would allow BA to repair project documentation and become a second
authority.

\subsubsection{Alternative 2 - localized
revalidation}\label{alternative-2---localized-revalidation}

\begin{verbatim}
changed revalidation context
    -> P0 STALE / PENDING_REVIEW
    -> source/context inconsistency localized
    -> DIAGNOSTIC_UNRESOLVED if no governed resolution exists
\end{verbatim}

After governance corrects/supersedes \texttt{FR-3.4}, the next baseline
can ground P1 from the corrected FR and apply T2 \texttt{REPLACE}
continuity to P0.

\begin{verbatim}
PASS
\end{verbatim}

This is the required behavior.

\subsection{\texorpdfstring{9. Why \texttt{revalidationContext} is not a
generic dependency
graph}{9. Why revalidationContext is not a generic dependency graph}}\label{why-revalidationcontext-is-not-a-generic-dependency-graph}

The pressure test tries three representations.

\subsubsection{Parentage only}\label{parentage-only}

Fails M4 because the influential Decision is sibling context.

\subsubsection{\texorpdfstring{BA2
\texttt{dependOn}}{BA2 dependOn}}\label{ba2-dependon}

Rejected because it would assert project-semantic prerequisite meaning
that may not exist.

\subsubsection{\texorpdfstring{Generic
\texttt{affects}}{Generic affects}}\label{generic-affects}

Rejected as too weak/underspecified for the current lower bound; it
invites a broad dependency ontology without saying what analytical
responsibility the edge carries.

\subsubsection{Retained minimum}\label{retained-minimum}

\begin{verbatim}
revalidationContext(targetElement, contextRef)
\end{verbatim}

with one precise analytical meaning: a material change to the context is
a reason to revalidate the target.

\subsection{10. Localized staleness algorithm
probe}\label{localized-staleness-algorithm-probe}

Starting changed governed source set \texttt{DeltaSource}, the lower
bound can seed review without scanning every BA element semantically
from scratch:

\begin{verbatim}
Seed S with elements whose:
  sourceLink intersects DeltaSource
  OR revalidationContext intersects DeltaSource

Repeat:
  if a derived element has a basisRef in S -> add it to S
  if an element has revalidationContext to a BA element in S -> add it to S
  if a proposition binds a referent
     replaced/retired through T2 review
     -> add it to S

For each element in S:
  mark PENDING_REVIEW / STALE for target baseline
  resolve semantically to RETAIN | REPLACE | RETIRE
\end{verbatim}

This is a responsibility-level algorithm sketch, not an implementation
data structure or scheduler.

\subsubsection{Negative control}\label{negative-control}

\begin{verbatim}
one changed file -> all BA elements stale
\end{verbatim}

remains rejected.

\subsection{11. Rule revision change
probe}\label{rule-revision-change-probe}

A derived BA element records the immutable rule revision used in its
historical origin.

If a later methodology revision changes the rule itself, historical
meaning is not rewritten. For a new materialization/review context the
element is re-evaluated under the explicitly selected current rule
revision.

Possible T2 outcomes:

\begin{verbatim}
same semantic result -> RETAIN
materially different result -> REPLACE
no longer applicable -> RETIRE or DIAGNOSTIC_UNRESOLVED
\end{verbatim}

This prevents methodological drift from silently changing the meaning
attributed to an older baseline.

\subsection{12. Feedback lineage
pressure}\label{feedback-lineage-pressure}

\subsubsection{Analysis-produced
candidate}\label{analysis-produced-candidate}

A downstream analysis may identify a gap in, for example, the
representation/FR relationship. To localize corrective feedback it can
carry:

\begin{verbatim}
BAElementRef@B0
GovernedSourceRef@B0
\end{verbatim}

obtained from Base Analysis provenance.

\subsubsection{Rejected correction}\label{rejected-correction}

If governance rejects the proposed correction:

\begin{verbatim}
analysis output exists
project documentation unchanged
BA source authority unchanged
\end{verbatim}

No new grounded BA element is introduced.

\subsubsection{Accepted correction}\label{accepted-correction}

If governance accepts the correction:

\begin{verbatim}
analysis/candidate
    motivates
new governed documentation B1
    grounds
BA@B1
\end{verbatim}

The BA@B1 \texttt{sourceLink} points to B1 governed documentation,
\textbf{not} to the finding/candidate that motivated it. T2 continuity
connects BA@B0 and BA@B1 where appropriate.

\subsubsection{Result}\label{result-4}

\begin{verbatim}
PASS WITHOUT NEW BA ANALYSIS IDENTITY
\end{verbatim}

The BA layer already exposes the handles required for downstream
analysis/change provenance. Exact AnalysisRecord/change-event lineage
remains downstream work and is not forced into BA3.

\subsection{13. Cross-layer authority
invariant}\label{cross-layer-authority-invariant}

The complete tested chain is:

\begin{verbatim}
governed docs B0
    -> BA B0
    -> downstream analysis
    -> corrective documentation candidate
    -> governed review
    -> governed docs B1
    -> BA B1
\end{verbatim}

The only transitions that establish project truth are governed-document
acceptance transitions. BA derivation rules and analysis outputs never
bypass that gate.

\subsection{14. Reopen checks}\label{reopen-checks}

\subsubsection{BA1}\label{ba1}

No new independently reusable project-semantic identity family is
required. Rule descriptors and revalidation links are metadata
responsibilities.

\begin{verbatim}
BA1 REOPEN: false
\end{verbatim}

\subsubsection{BA2}\label{ba2}

No new project-semantic operator is required. In particular,
\texttt{revalidationContext} is not \texttt{dependOn} and does not enter
the BA2 operator registry.

\begin{verbatim}
BA2 REOPEN: false
\end{verbatim}

\subsubsection{BA3-T1}\label{ba3-t1}

T1 is refined by role-binding \texttt{derivationBasis}; its
provenance/origin responsibilities remain valid.

\begin{verbatim}
T1 REOPEN: false
\end{verbatim}

\subsubsection{BA3-T2}\label{ba3-t2}

T2 staleness/lifecycle semantics remain valid; T3 supplies the
previously open effective-context representation needed to seed them.

\begin{verbatim}
T2 REOPEN: false
\end{verbatim}

\subsection{15. Disposition table}\label{disposition-table}

\begin{verbatim}
Plain untyped derivationBasis list                 REJECTED
Role-bound derivationBasisBinding                  REQUIRED
Mutable/opaque derivation rule                     REJECTED
Immutable inspectable rule revision                REQUIRED
Universal executable rule language                 REJECTED
Semantic replay reproducibility                    REQUIRED
M4 parentage-only change detection                 REJECTED
Explicit revalidationContext                       REQUIRED
revalidationContext -> BA2 dependOn                REJECTED
Global BA invalidation on any source change        REJECTED
Localized staleness propagation                    PASS
Auto-repair BA from sibling source                 REJECTED
Provider mapping hidden in tool                    REJECTED
Governed provider mapping as derivation basis      REQUIRED
  when the mapping itself is derived
Analysis finding/candidate as BA source             REJECTED
New BAE family                                     NOT FORCED
BA1/BA2 reopen                                     NOT TRIGGERED
Distinct BA3 closure review                        REQUIRED
\end{verbatim}

\subsection{16. Exit condition
assessment}\label{exit-condition-assessment}

BA3-T3 exits with:

\begin{verbatim}
COMPLETED / PROVISIONAL PASS WITH DERIVATION-IMPACT REFINEMENT
\end{verbatim}

The T1/T2 lower bounds survive with two refinements:

\begin{enumerate}
\def\labelenumi{\arabic{enumi}.}
\tightlist
\item
  derivation basis is role-bound and references an immutable inspectable
  rule revision;
\item
  explicit \texttt{revalidationContext} captures
  non-source/non-derivation governed context needed to seed localized
  review.
\end{enumerate}

The feedback loop requires no new BA authority or BAE family.

A distinct BA3 closure review is still required because T1-T3 have
accumulated multiple orthogonal metadata responsibilities that must be
adversarially checked for redundancy, contradiction and cross-corpus
sufficiency before BA4 begins.

\subsection{17. Next authorized
microstep}\label{next-authorized-microstep}

Only after this package is reviewed, committed, pushed and remotely
verified:

\begin{quote}
\textbf{BA3-T4 - provenance, identity/lifecycle and
derivation/change-impact closure review.}
\end{quote}

Do not start BA4 until that review either closes BA3 or reopens the
smallest failed BA3 responsibility.

\end{document}
